\textbf{结论名称：}极值点偏移：对数均值不等式法（标准路径）

\textbf{适用条件：} 函数 \(f(x)\) 满足 \(f(x_1)=f(x_2)\) 且可变形出 \(\ln x_1-\ln x_2\) 与 \(x_1-x_2\) 的比值结构；常用于含 \(\ln x\) 或可化为此类的函数，证明两根之和（或积）相对于极值点的不等式。

\textbf{变量说明：} 设 \(x_1,x_2\) 是 \(f(x)\) 的两个不等实根，且 \(x_1,x_2>0\)；\(x_0\) 为 \(f(x)\) 的极值点。

\textbf{核心公式：}
\[
\boxed{\sqrt{x_1x_2} < \frac{x_1-x_2}{\ln x_1-\ln x_2} < \frac{x_1+x_2}{2}\qquad (x_1,x_2>0,\;x_1\neq x_2)}
\]
常用上界形式：\(\displaystyle\frac{x_1-x_2}{\ln x_1-\ln x_2}<\frac{x_1+x_2}{2}\)，等价变形为 \(\displaystyle\ln\frac{x_1}{x_2}>\frac{2(x_1-x_2)}{x_1+x_2}\)。

\textbf{使用场景：} 处理含对数结构的极值点偏移问题。将 \(f(x_1)=f(x_2)\) 整理成 \(\displaystyle\frac{x_1-x_2}{\ln x_1-\ln x_2}=k\)（或相关形式），代入对数均值不等式，消去超越项，直接得到 \(x_1+x_2\) 或 \(x_1x_2\) 与常数的不等关系，是证明 \(x_1+x_2>2x_0\) 或 \(x_1x_2<x_0^2\) 的标准化简路径。

\textbf{简单例子：} 已知 \(f(x)=x\mathrm{e}^{-x}\)，\(f(x_1)=f(x_2)\) 且 \(x_1\neq x_2\)，求证 \(x_1+x_2>2\)。
由 \(x_1\mathrm{e}^{-x_1}=x_2\mathrm{e}^{-x_2}\) 两边取对数得 \(\ln x_1-x_1=\ln x_2-x_2\)，整理有
\[
\frac{x_1-x_2}{\ln x_1-\ln x_2}=1.
\]
不妨设 \(x_1>x_2>0\)，代入对数均值不等式上界：
\[
1=\frac{x_1-x_2}{\ln x_1-\ln x_2}<\frac{x_1+x_2}{2}\implies x_1+x_2>2.
\]